\documentclass{auxiliares/umemoria} 
\depto{Departamento de Matemática y Ciencia de la Computación}
\author{Nombres y apellidos}
\title{Linea 1{\vspace{1.5ex}\break}Linea 2}
%si el titulo supera una linea agrgar {\vspace{1.5ex}\break} al final de donde se desea la primera linea para que se ajuste el texto
\bajadatitulo{Primer Avance}

\titulopdf{Nombre pdf}

% Grado al que opta cambiar para masculino y femenino
\tesis{Ingeniero Estadístico}

% Profesor guía 
\guia{Nombres y apellidos}

% Miembros comisión
\comisiona{Este aparece en primer lugar} % Este aparece en primer lugar
\comisionb{Este aparece en segundo lugar} % Este aparece en segundo lugar

% Año en que se da el examen de título/grado (defensa) o 
% ultima matricula de pregrado
\anho{AÑO}

% Lugar donde se encuentra la bibliografia
\addbibresource{g_bibliografia.bib}

% Hints
% formulas que se repiten como notaciones de vectores que son muy extensas de escribir en cada apartado
\input{auxiliares/hints} 

    %PÁGINAS PRELIMINARES
%%%%%%%%%%%%%%%%%%%%%%%%%%%%%%%%%%
%(1) PORTADA 			
%(2) DERECHOS DE AUTOR 		
%(3) RESUMEN Y ABSTRACT 	
%(4) DEDICATORIA(OPTATIVA)
%(5) AGRADECIMIENTOS(OPTATIVA)	
%(6) TABLA DE CONTENIDOS
%(7) ÍNDICE DE TABLAS	
%(8) ÍNDICE DE ILUSTRACIONES
%%%%%%%%%%%%%%%%%%%%%%%%%%%%%%%%%%

\begin{document}
\frontmatter
\maketitle

% Derechos de autor / omitir cuando se envia a biblioteca
\begin{derechoautor}
\copyright \,Nombres Apellidos, Año.\\
\textbf{Todos los derechos reservados. Chile.}
\end{derechoautor}

% Resumen breve
\begin{resumen}
\lipsum[1-2]\\

\noindent\textbf{Palabras claves:}  % cada palabra comienza con mayuscula separadas por coma.
\end{resumen}

% Abstract
\begin{abstract}
\lipsum[1-2]\\

\noindent\textbf{Keywords:}  % cada palabra comienza con mayuscula separadas por coma.
\end{abstract}

% Dedicatoria corta o frase inpiradora, poema, etc.
\begin{dedicatoria}
``Frase, poema o lo que quieras poner de manera breve.''\\
\vspace*{0.4cm}
- Persona que lo dijo\\
\vspace*{2.6cm}
``Dedicatoria a familia''
\end{dedicatoria}

% Agradecimientos
\begin{thanks}
\input{i_agradecimientos}
\end{thanks}

\cleardoublepage
\phantomsection
\tableofcontents % índice general
\listoftables % índice de tablas
\listoffigures % índice de figuras

\mainmatter

\input{a_introduccion}
\input{b_obj_metodologia}
\input{c_marco_teorico}
\input{d_resultados}
\input{e_discusion}
\input{f_conclusion}

% Bibliografia
\begingroup
    \setlength\bibitemsep{1.5ex}
    \setstretch{1}
    \cleardoublepage % hace lo mismo que el clearpage pero no altera los contadores
    \phantomsection % para ajustar las sectiones fantasmas
    {\addcontentsline{toc}{chapter}{Referencias bibliográficas}\printbibliography[title={Referencias bibliográficas}]}
\endgroup

% opcional ...
% Para setear el capitulo del anexo 
\titleformat{\chapter}[display]{\bf\huge}{\chaptertitlename$\hspace{2.7mm}$\thechapter}{10pt}{}
\begin{appendices}
 \input{h_anexo}
\end{appendices}
\end{document}